% phi=.5 x^2,base x=-1,p=2; secants at theta1,.5,.25 approach tangent slope-1.
\begin{tikzpicture}[font=\small,>={Latex[length=2mm]},x=1.8cm,y=1.8cm]
 \draw[->,black] (-1.6,0)--(1.6,0) node[right] {$s$};
 \draw[->,black] (0,-.65)--(0,1.45) node[above] {$\phi$};
 \draw[thick,domain=-1.45:1.5,samples=60] plot(\x,{.5*\x*\x});
 \draw[dashed,black] (-1,.5)--(1,.5);
 \draw[blue!70!black,dash dot] (-1,.5)--(0,0);
 \draw[blue!70!black,densely dotted,thick] (-1,.5)--(-.5,.125);
 \draw[thick,domain=-1.4:.1] plot(\x,{-\x-.5});
 \fill (-1,.5) circle(1.5pt) node[above left] {$x$};
 \fill (1,.5) circle(1.5pt) node[above right] {$x+p$};
 \fill (0,0) circle(1.3pt);\fill (-.5,.125) circle(1.3pt);
 \draw[->,thick] (-.05,.1) to[bend right=20] (-.88,.52);
 \node[anchor=east] at (-.6,.9) {$\theta\downarrow0$};
 \node[anchor=west,align=left,text width=62mm] at (2.1,.55)
  {$x+\theta p$ approaches $x$.\\[6pt]
   Secants approach the tangent\\at the base point.\\[6pt]
   Subtract $\phi(x)$, divide by\\positive $\theta$, then take the limit.};
\end{tikzpicture}
